\chapter{Introduction}
\label{ch:introduction}

India is the world's third-largest consumer of crude oil and imports roughly 85 per cent of its requirements. This dependence means that movements in global oil prices feed directly into domestic costs, affecting everything from transport and manufacturing to household cooking fuel. When global oil prices rise, the cost of imported crude increases India's import bill, puts downward pressure on the rupee, and raises input costs across sectors. The resulting inflationary pressure concerns policymakers at the Reserve Bank of India (RBI), which targets a 4 per cent headline CPI inflation rate under its flexible inflation targeting framework.

A natural question follows: does a fall in oil prices bring inflation down by just as much as a rise pushes it up? The answer matters for monetary policy, fiscal planning, and welfare. If oil price increases pass through to consumer prices more strongly than decreases, a pattern the literature calls ``rockets and feathers,'' then the average inflationary effect of oil price volatility is positive, even when oil prices fluctuate symmetrically around their mean. Households would bear a net inflationary cost from oil price instability, and the central bank would need to account for this asymmetry in its inflation forecasts.

\section{Research Question}

This dissertation examines the short-run pass-through of global oil price shocks to India's headline CPI inflation over the period April 2004 to December 2024. Two specific questions guide the analysis:

\begin{enumerate}[label=(\roman*)]
\item Do positive oil price shocks raise monthly CPI inflation by an economically meaningful amount?
\item Is the pass-through from oil price increases larger than that from oil price decreases?
\end{enumerate}

The first question asks whether oil prices matter for Indian inflation in the short run. The second asks whether the transmission is asymmetric.

\section{Motivation}

Several features of India's economy make this question worth studying in detail.

First, India's fuel pricing regime has changed considerably during the sample period. Petrol prices were formally deregulated in June 2010, and diesel prices followed in October 2014. Before deregulation, the government absorbed much of the oil price volatility through subsidies, which may have dampened pass-through to consumer prices. After deregulation, domestic retail fuel prices began tracking international prices more closely, at least in principle. This institutional change provides a natural setting for comparing pass-through across two distinct policy regimes.

Second, the exchange rate adds a layer of complexity that is specific to India's situation. India pays for crude oil in US dollars, so a depreciation of the rupee against the dollar amplifies the domestic-currency cost of any given oil price increase. This means that India-specific inflationary episodes can arise even during periods of stable global oil prices, if the rupee weakens. Separating the oil price channel from the exchange rate channel requires explicit attention, which the robustness analysis in this dissertation provides.

Third, India's CPI basket assigns roughly 46 per cent weight to food items and only about 7 per cent directly to fuel and light. This composition means that headline CPI is a noisy measure of energy cost pass-through. Oil prices affect headline CPI through multiple channels (direct fuel costs, transport costs embedded in food distribution, and manufacturing input costs), but the signal gets diluted by weather-driven food inflation and other shocks that have nothing to do with oil. This makes it difficult to obtain precise estimates of oil pass-through in headline CPI, a difficulty that should be acknowledged upfront rather than glossed over.

\section{Approach}

The main empirical model is an asymmetric autoregressive distributed lag (ADL) specification in first differences. Oil price changes are split into positive and negative components, and the cumulative effects of each are estimated separately. The oil lag length is fixed at three months, and the autoregressive CPI lag order is selected by the Akaike information criterion over one to four lags on a common estimation sample. All hypothesis tests use Newey--West heteroskedasticity and autocorrelation consistent (HAC) standard errors, so the inference is robust to serial correlation and heteroskedasticity in the residuals.

For robustness, a separate specification replaces the composite rupee oil price with Brent crude in US dollars and includes the exchange rate as a separate regressor. This allows the world oil price effect and the exchange rate effect to be estimated independently. An appendix extends the analysis to the official CPI Fuel and Light sub-index, which is more directly exposed to energy cost movements than headline CPI.

\section{Structure of the Dissertation}

Chapter 2 reviews the relevant literature on oil price shocks, inflation, and asymmetric pass-through. Chapter 3 describes the data and econometric methodology. Chapter 4 presents the main empirical results, including unit root tests, the baseline symmetric model, the main asymmetric model, and sub-sample comparisons. Chapter 5 covers robustness checks. Chapter 6 discusses the findings and their policy relevance. Chapter 7 concludes.
